\documentclass[10pt,a4paper]{amsart}
\usepackage[margin=19mm]{geometry}
\usepackage{amsmath,amssymb,mathtools}
\usepackage[scr=rsfs,cal=euler]{mathalfa}
\usepackage{xfrac}
\usepackage[unicode,hidelinks,bookmarks=false]{hyperref}
\newcommand{\ie}{i.e.\ }
\newcommand{\cf}{cf.\ }
\newcommand{\resp}{respectively\ }
\newcommand{\Subsection}{\subsection}
\providecommand{\nobreakdash}{\nobreak\hbox{-}}
\setlength{\emergencystretch}{2em}
\multlinegap0pt
\usepackage{amssymb}
\usepackage{float}
\usepackage{bm}
\usepackage{tikz}
\newtheorem{lemma}{Lemma}
\newtheorem{claim}{Claim}
\def\qed{\hfill$\Box$\vspace{12pt}}
\title{Signless Laplacian spectral conditions for rainbow matchings in a collection of bipartite graphs}
\author{Zhiwei Guo, Nanxi Pan, Li Li, Yangyang Chen}
\date{Source: arXiv:2608.26712v2 (CC BY 4.0)}
\begin{document}
\maketitle

\noindent\emph{Source and licence.} Signless Laplacian spectral conditions for rainbow matchings in a collection of bipartite graphs. Version arXiv:2608.26712v2 (CC BY 4.0), distributed by arXiv under the Creative Commons Attribution 4.0 International licence, http://creativecommons.org/licenses/by/4.0/, as recorded in the arXiv licence metadata for this exact version and shown on the abstract page https://arxiv.org/abs/2608.26712v2. Full licence notice: \texttt{licenses/arxiv-2608.26712-CC-BY-4.0.txt}.\par\smallskip

\noindent\textit{Editorial scope.} Counted unit: the article's claim claim1a (source0.tex lines 293--295). The article's complete original proof of it is delivered with it (source0.tex lines 297--323, one \texttt{proof} environment of 27 lines). No mathematics is written, repaired or re-derived: the statement and the proof are the article's own lines, in the article's own notation, and no retained line is substituted or re-flowed.  Retained uncounted as the source-local context the proof needs: lines 212--214; lines 253--255; lines 259--261; lines 285--288.  Nothing was omitted from any retained span, so the package carries no omission record.  The retained lines cite 2 works, whose entries are transcribed verbatim from the article's own bibliography source in the e-print and delivered with short editorial labels recorded in the item's reference mapping.  No retained line contains a figure environment, an image-inclusion command or drawing code, so no image is delivered and the package declares no asset dependency.  Editorial formatting only, in addition: an AMS \texttt{amsart} preamble that restates the article's own declarations and macros used by the retained lines, namely the environments \texttt{lemma}, \texttt{claim} on separate counters, together with the standard packages those lines need.  The retained environments print in this document as Lemma 1, Claim 1 in order of appearance, whereas the article prints the same items with the numbering of its own full document.  The complete native source of the article as distributed in the arXiv e-print for this exact version is bundled unchanged as \texttt{sources/208657/source0.tex}.
\begin{lemma}\label{lemma3b}
Let $G$ be a shifted graph. Then $\{{{y_1},{y_2}}\}\in E(G)$ implies $\{ {{x_1},{x_2}}\} \in E(G)$ for each $\{{{x_1},{x_2}}\},\{{{y_1},{y_2}}\} \subseteq [n]$ such that ${x_i}\le {y_i}$ for $i \in \{1,2\}$.
\end{lemma}

\begin{lemma}[Godsil and Royle \cite{Godsil}]\label{lemma5b}
Let $G$ be a graph. If $\pi$ is an equitable partition of $V(G)$ corresponding to $Q(G)$, then $\lambda_{\max}(Q(G)/\pi)=q(G)$.
\end{lemma}

\begin{lemma}[Brouwer and Haemers \cite{Brouwer}]\label{lemma6b}
Let $H$ be a subgraph of a connected graph $G$. Then $q(H)\le q(G)$ and the equality holds if and only if $H \cong G$.
\end{lemma}

\begin{equation}
\label{eq:shift-lower}
q(S_i)\ge q(G_i)\ge b+k-1
\end{equation}

\begin{claim}\label{claim1a}
For each $i\in[k]$, the graph $S_i$ contains $e_2,\ldots,e_{k-1}$.
\end{claim}

\begin{proof}
For the case of $k=2$, the assertion is vacuous. We now assume that $k\ge3$. Suppose that $e_r\notin E(S_i)$ for some $i\in[k]$ and $2\le r\le k-1$. If $uv'\in E(S_i)$ with $u\ge r$ and $v'\ge k-r+1$, then by Lemma~\ref{lemma3b}, we deduce that $e_r=r(k-r+1)'\in E(S_i)$, a contradiction. Hence, no edge of $S_i$ joins $\{r,\ldots,a\}$ to $\{(k-r+1)',\ldots,b'\}$.

Let $H_r$ be the graph obtained from $K_{a,b}$ by deleting all edges between these two sets. Since $r-1\ge1$ and $k-r\ge1$, the vertices $1,\ldots,r-1$ are adjacent to all of $Y$. However, the vertices $1',\ldots,(k-r)'$ are adjacent to all of $X$. It follows that $H_r$ is connected. Thus, $S_i$ is a spanning subgraph of $H_r$. By Lemma~\ref{lemma6b}, we have $q(S_i)\le q(H_r)$. Let $X_1=[r-1]$, $X_2=[r,a]$, $Y_1=[1,k-r]'$ and $Y_2=[k-r+1,b]'$. The partition $\pi=\{X_1,X_2,Y_1,Y_2\}$
is equitable with respect to $Q(H_r)$. Indeed, the vertices in \(X_1,X_2,Y_1,Y_2\) have degrees \(b,k-r,a,r-1\), respectively, and their neighbors are distributed among these four sets according to the quotient matrix
\[B=\begin{pmatrix}
b&0&k-r&b-k+r\\
0&k-r&k-r&0\\
r-1&a-r+1&a&0\\
r-1&0&0&r-1
\end{pmatrix},
\]
where the entries of \(B\) are obtained by summing the corresponding rows of \(Q(H_r)\) over the four parts.

Put $p=k-r$, $h=r-1$, and $s=p+h=k-1$, and $d=b-a\ge0$. The characteristic polynomial of $B$ is
$g(x):=\det(xI-B)=x^4-(a+b+t)x^3+(2ph+pb+ah+ab)x^2-ph(a+b)x$. By Lemma~\ref{lemma5b}, $q(H_r)$ is the largest eigenvalue of $B$, and hence $q(H_r)$ is a root of $g$. Substituting $a=b-d$, $t=p+h$, and $x=b+t+y$ gives
$g(b+t+y)=y^4+c_3y^3+c_2y^2+c_1y+c_0$, where
\begin{align*}
c_3={}&2b+d+3h+3p,\\
c_2={}&b^2+2bd+4bh+4bp+2dh+3dp+3h^2+8hp+3p^2,\\
c_1={}&b^2d+b^2h+b^2p+2bdh+4bdp+2bh^2+6bhp+2bp^2\\
&\quad+dh^2+5dhp+3dp^2+h^3+7h^2p+7hp^2+p^3,\\
c_0={}&p(b+h+p)\bigl(bd+2dh+dp+2h^2+2hp\bigr).
\end{align*}
Since $p,h\ge1$, $b>0$, and $d\ge0$, each of the four coefficients is positive. The positivity of these coefficients implies $g(b+s+y)>0$ for every $y\ge0$. Hence, $g(x)>0$ whenever $x\ge b+s$. Since $q(H_r)$ is a root of $g$, and $g(x)>0$ for all
$x\ge b+k-1$, we deduce that $q(H_r)<b+k-1$, contrary to Inequality~\eqref{eq:shift-lower}.\qed
\end{proof}

\begin{thebibliography}{99}
\bibitem[Brouwe]{Brouwer} A.E. Brouwer, W.H. Haemers, Spectra of graphs, Universitext, Springer, New York, 2012.
\bibitem[Godsil]{Godsil} C. Godsil, G.F. Royle, Algebraic Graph Theory, Springer, New York, 2001.
\end{thebibliography}
\end{document}
